\chapter{物理因素及其对健康的影响}
\setcounter{chapter}{5}
\chapterintro{
本章为课程重要章节（15学时）。

\textbf{【教学大纲要求】}

\imp{掌握：}生产环境中常见的物理因素及物理因素的特点\keyword{*}；高温作业类型；高温作业对机体生理功能的影响\keyword{*}（体温调节、水盐代谢、循环系统、消化系统、泌尿系统、神经系统）；热适应的生理学意义；中暑（致病因素、发病机制与分型\keyword{*}、诊断、急救与治疗）；防暑降温措施（技术措施、保健措施与组织措施）；我国高温作业的卫生标准；高、低气压作业特点及对机体的影响；减压病与高原病；生产性噪声的来源及分类；声强与声强级、声压与声压级、频谱、响度、等响曲线、计权声级\keyword{*}；噪声对人体的影响（听觉系统、神经系统、心血管系统等）；噪声卫生标准；振动的概念、分类与接触机会；振动的物理参数\keyword{*}（频率、位移、速度、加速度、振动频谱、共振频率、振动轴向）；全身振动与局部振动对健康的影响；手臂振动病；非电离辐射（射频辐射的生物学效应、红外辐射、紫外辐射、激光对健康的影响）；电离辐射的卫生防护。

\imp{熟悉：}低温作业；影响噪声对机体作用的因素；防止噪声危害的措施；影响振动对机体作用的因素；预防振动职业危害的措施；电离辐射的基本概念、接触机会、作用及其影响因素、临床表现。

（注：\keyword{*} 标示教学难点）
}

% ----- 5.1 高温作业 -----
\section{高温作业与中暑}

\subsection{高温作业的基本概念}

\begin{defbox}
\textbf{高温作业}：工作地点有高气温，或有强烈的热辐射，或伴有高气湿相结合的异常气象条件，WBGT指数超过规定限值的作业。

\textbf{高温作业类型：}
\begin{enumerate}[leftmargin=*]
    \item \imp{高温、强热辐射作业（干热型）}：炼钢、炼铁、铸造、玻璃熔制等，特点是气温高、热辐射强、相对湿度低
    \item \imp{高温、高湿作业（湿热型）}：纺织、印染、造纸、深井煤矿等，特点是气温高、湿度大（>80\%）、热辐射不强
    \item \imp{夏季露天作业}：农业、建筑、搬运等，受太阳直接辐射和地面的二次辐射
\end{enumerate}
\end{defbox}

\subsection{高温作业对机体的影响*}

\begin{keybox}
\textbf{（一）体温调节与水盐代谢}
\begin{itemize}[leftmargin=*]
    \item 热平衡公式：$S = M - E \pm R \pm C_1 \pm C_2$（S：热蓄积，M：代谢产热，E：蒸发散热，R：辐射获/散热，C$_1$：对流获/散热，C$_2$：传导获/散热）
    \item 高温下机体主要通过\keyword{出汗蒸发}散热
    \item 当环境温度超过皮肤温度（35$^\circ$C）时，辐射和对流散热停止，出汗蒸发成为唯一散热方式
    \item 大量出汗导致水盐丢失（日出汗量可达6--8L），若不及时补水补盐，可引起水盐代谢紊乱
\end{itemize}

\textbf{（二）循环系统}
\begin{itemize}[leftmargin=*]
    \item 心率加快（每升高1$^\circ$C，心率增加10次/分）、血压先升后降
    \item 外周血管扩张、内脏血管收缩——血液重新分配
    \item 心脏负担加重，长期高温作业可致心肌损害
\end{itemize}

\textbf{（三）消化系统}
\begin{itemize}[leftmargin=*]
    \item 消化液分泌减少、消化酶活性降低→食欲减退、消化不良
    \item 胃肠运动减弱
\end{itemize}

\textbf{（四）泌尿系统}
\begin{itemize}[leftmargin=*]
    \item 肾血流量减少、尿液浓缩→肾负担加重
    \item 可出现蛋白尿、管型尿
\end{itemize}

\textbf{（五）神经系统}
\begin{itemize}[leftmargin=*]
    \item 中枢神经系统受抑制——注意力不集中、反应迟钝、肌肉工作能力下降、作业能力降低
\end{itemize}
\end{keybox}

\subsection{热适应与热习服}

\begin{defbox}
\textbf{热适应（Heat Adaptation）}：机体在长期高温环境下产生的形态、结构和生理功能的适应性变化。热适应是机体与环境之间形成的相对稳定的适应关系，涉及体温调节、水盐代谢、心血管功能等多系统的综合性改变。

\textbf{热适应与热习服的区别：}
\begin{table}[H]
\centering
\caption{热适应与热习服的区别}
\begin{tabularx}{\textwidth}{lXX}
\hline
\textbf{特征} & \textbf{热适应（Heat Adaptation）} & \textbf{热习服（Heat Acclimatization）} \\
\hline
形成时间 & 长期（世代/常年） & 短期（数天至数周，一般7--14天）\\
机制 & 基因表达和生理结构改变 & 生理功能补偿性调整 \\
可逆性 & 较稳定，不易消退 & 脱离高温后数天至数周消退 \\
接触条件 & 多年在高温环境下生活和工作 & 短期反复接触高温 \\
出汗特点 & 出汗效率高，NaCl损失少 & 出汗量增加，NaCl损失减少 \\
循环系统 & 心率较稳定，心血管负荷小 & 心率适应性降低 \\
\hline
\end{tabularx}
\end{table}

\textbf{热习服的机制：}
\begin{enumerate}[leftmargin=*]
    \item 体温调节改善：出汗阈值降低，出汗量增加，散热效率提高
    \item 水盐代谢调整：醛固酮分泌增加减少钠盐丢失
    \item 心血管系统适应：心率降低，每搏输出量增加
    \item 细胞热休克蛋白（HSP）表达增强，提高细胞耐热能力
\end{enumerate}

\textbf{热习服的临床意义：}热习服可显著提高高温作业能力和预防中暑。缺乏热习服者（如新工、休假返工者）是中暑的高危人群。高温作业前应安排3--7天的热习服期。
\end{defbox}

\subsection{中暑（Heat Illness）}

\begin{keybox}
\textbf{中暑的临床分型*：}

\textbf{1. 热射病（Heat Stroke）*}：
\begin{itemize}[leftmargin=*]
    \item 高温高湿、强热辐射；体内蓄热，体温调节紊乱
    \item 特征：\keyword{高热（体温>40$^\circ$C）+ 汗闭 + 意识障碍（昏迷）}
    \item 机制：体温调节中枢功能衰竭→产热>散热
\end{itemize}

\textbf{2. 日射病（Sun Stroke）}：
\begin{itemize}[leftmargin=*]
    \item 夏季露天作业；颅内受热导致脑膜和脑组织充血
    \item 特征：头（颅）温升高、有汗、恶心呕吐、昏迷
\end{itemize}

\textbf{3. 热痉挛（Heat Cramps）}：
\begin{itemize}[leftmargin=*]
    \item 高温、强热辐射；大量出汗后仅补水未补盐→低钠血症→肌肉痉挛（腓肠肌最多见）
    \item 体温多正常或轻度升高，无昏迷
\end{itemize}

\textbf{4. 热衰竭（Heat Exhaustion）}：
\begin{itemize}[leftmargin=*]
    \item 高温、强热辐射；心血管功能失代偿→有效循环血量不足→虚脱
    \item 特征：冷汗、面色苍白、皮肤湿冷、脉搏细弱、血压降低、昏厥、体温升高
\end{itemize}

\textbf{中暑各型的临床特点对比：}

\begin{table}[H]
\centering
\caption{中暑四种类型的临床特点对比}
\begin{tabularx}{\textwidth}{lXXXX}
\hline
\textbf{特征} & \textbf{热射病} & \textbf{日射病} & \textbf{热痉挛} & \textbf{热衰竭} \\
\hline
出汗 & 汗闭 & 有汗 & 大量出汗 & 冷汗 \\
体温 & 体温↑ & 头（颅）温↑ & 体温不高 & 体温↑ \\
意识 & 昏迷 & 昏迷 & 无昏迷 & 昏迷 \\
其他 & CNS症状 & 恶心呕吐 & 肌肉痉挛 & 晕厥 \\
\hline
\end{tabularx}
\end{table}

\textbf{中暑的诊断（GBZ 41-2019）：}
\begin{itemize}[leftmargin=*]
    \item 中暑先兆：头昏/痛、口渴、多汗、疲乏、心悸、注意力不集中、动作不协调，体温正常或稍高
    \item 重症中暑：热射病（含日射病）、热痉挛、热衰竭
\end{itemize}

\textbf{救治原则：}
\begin{enumerate}[leftmargin=*]
    \item 立即脱离高温环境，恢复热平衡和水盐平衡、防止休克
    \item \imp{物理降温}——冷水/冰水/酒精反复擦浴、扇风（注意：肛温低于38$^\circ$C时应立即停止降温）
    \item \imp{药物降温}：影响体温调节中枢减少产热；扩张周围血管加速散热；松弛肌肉减少肌肉震颤；降低细胞氧消耗使机体更好地耐受缺氧
    \item 补充水分和电解质（口服补盐液、静脉输液）
    \item 对症治疗（呼吸循环支持、防治脑水肿等）
\end{enumerate}

\textbf{预防措施：}
\begin{enumerate}[leftmargin=*]
    \item 技术措施：合理设计、改革工艺；隔热（水箱、瀑布水幕、石棉、炉渣、耐火砖等）；通风（自然通风+机械通风）
    \item 保健措施：合理供应含盐饮料；合理膳食（补充蛋白质、维生素）；热适应训练
    \item 组织措施：合理安排作息制度（轮换作业、延长休息）、高温作业标准（WBGT）；医学监测（职业禁忌证）
\end{enumerate}
\end{keybox}

\subsection{我国高温作业卫生标准}

以\keyword{湿球黑球温度指数（WBGT）}为评价指标。WBGT采用了自然湿球温度（$t_{nw}$）、黑球温度（$t_g$）和干球温度（$t_a$）三种参数：
\begin{itemize}[leftmargin=*]
    \item 室内作业：WBGT = $0.7t_{nw} + 0.3t_g$
    \item 室外作业：WBGT = $0.7t_{nw} + 0.2t_g + 0.1t_a$
\end{itemize}
WBGT综合反映气温、气湿、气流和热辐射四种物理因素对机体的作用。根据体力劳动强度和接触时间制定不同限值。

% ----- 5.2 异常气压 -----
\section{异常气压}

\subsection{高气压作业}

\begin{defbox}
\textbf{高气压作业}：潜水作业（每下沉10.3m增加1个大气压）和潜涵作业（隧道、桥墩等地下工程）。

\textbf{高气压对机体的影响：}
\begin{itemize}[leftmargin=*]
    \item 减压过程：溶解在体内的氮气（N$_2$）从组织中释放形成气泡
    \item 若减压速度过快→\keyword{减压病（Decompression Sickness）}
\end{itemize}
\end{defbox}

\begin{keybox}
\textbf{减压病的发病机制*：}在高压环境下，氮气大量溶解于组织和体液中。若减压过快，溶解的氮气迅速释放形成气泡，栓塞血管和压迫组织→组织缺血缺氧。

\textbf{临床表现：}
\begin{itemize}[leftmargin=*]
    \item 皮肤瘙痒、大理石样斑纹
    \item 关节和肌肉剧痛（\textbf{屈肢症}——被迫屈曲患肢以缓解疼痛）
    \item 脊髓损害——截瘫、感觉障碍
    \item 脑部气泡栓塞——昏迷
\end{itemize}

\textbf{治疗：}及时送入\keyword{加压舱}进行\keyword{加压治疗}（唯一有效治疗）+吸氧+对症

\textbf{预防：}严格遵守减压程序（分段减压、足够减压时间）；控制作业时间；健康监护（禁忌证：呼吸系统疾病、心血管疾病、肥胖）
\end{keybox}

\subsection{低气压作业}

高原作业（海拔>3000m）。主要问题是\keyword{低氧}——海拔越高，氧分压越低。机体通过增加肺通气量、心率加快、红细胞增多等进行代偿适应。急性高原病：头痛、恶心、呼吸困难、肺水肿、脑水肿。

% ----- 5.3 噪声 -----
\section{生产性噪声}

\subsection{基本概念}

\begin{defbox}
\textbf{生产性噪声}：在生产过程中产生的，频率和强度无规律组合、使人感到厌烦或影响健康的声音。

\textbf{噪声的分类（按来源）：}
\begin{enumerate}[leftmargin=*]
    \item 机械性噪声：机械撞击、摩擦、转动产生
    \item 流体动力性噪声：气体/液体流速或压力突变产生
    \item 电磁性噪声：交变电磁场产生
\end{enumerate}

\textbf{噪声的分类（按频谱）：}
\begin{itemize}[leftmargin=*]
    \item 低频噪声（<300Hz）、中频噪声（300--800Hz）、高频噪声（>800Hz）
\end{itemize}

\textbf{噪声的分类（按时间分布）：}
\begin{itemize}[leftmargin=*]
    \item 稳态噪声：声压级波动<5dB
    \item 脉冲噪声：声压级波动$\geq$5dB，持续时间短
    \item 非稳态噪声：声压级波动$\geq$5dB
\end{itemize}
\end{defbox}

\subsection{噪声的物理特性与评价指标}

\begin{keybox}
\textbf{声强与声强级：}声强——单位时间内通过单位面积波阵面的声能。人耳能感受的声强范围极宽（10$^{-12}$--1 W/m$^2$）。

\textbf{声压与声压级（SPL）：}$L_p = 20 \lg \frac{p}{p_0}$ (dB)

\textbf{响度级与等响曲线：}以1000 Hz纯音为基准，将不同频率的声音与基准音比较得出的主观感觉量，单位为\keyword{方（phon）}。

\textbf{A计权声级（A声级）*}：用A计权网络测量的声压级，模拟人耳对40方纯音的响应，单位为\keyword{dB(A)}。是评价环境噪声和工业噪声最常用的指标。

\textbf{等效连续A声级（$L_{Aeq}$）}：在规定时间内，某一连续稳态声的A声级具有与实际噪声相同的能量。用于评价非稳态噪声。

\textbf{累积百分声级（$L_N$）}：在规定的测量时间内，有N\%的时间超过的A声级。如$L_{10}$、$L_{50}$、$L_{90}$。
\end{keybox}

\subsection{噪声对听觉系统的影响}

\begin{keybox}
\textbf{听觉适应}：短时间暴露于强噪声环境中，听力下降（听阈提高10--15dB），脱离后数分钟恢复。属于生理性保护反应。

\textbf{听觉疲劳}：较长时间暴露于强噪声，听力明显下降（听阈提高15--30dB），脱离后数小时至十几小时才能恢复。属生理性反应的延续。

\textbf{噪声性听力损伤（永久性听阈位移，NIPTS）*}：
\begin{itemize}[leftmargin=*]
    \item 长期暴露于强噪声引起的不可逆性听力损害
    \item \imp{高频听力下降}（3000--6000 Hz，尤其4000 Hz出现V形凹陷——\keyword{4kHz听谷}）为最早特征
    \item 随着损害加重，语言频率（500--2000 Hz）听力也下降
    \item 听力曲线由高频向低频扩展
\end{itemize}

\textbf{噪声聋（Noise-Induced Deafness）}：长期接触生产性噪声引起的以高频听力下降为主、继而语言频率听力下降的\keyword{感音性耳聋}。属于法定职业病。诊断依据GBZ49-2014，噪声作业接触史需\keyword{3年以上}。

\textbf{爆震性耳聋}：一次或数次接触高强度脉冲噪声（如爆破、枪炮发射）引起的急性听觉器官损伤。常伴鼓膜破裂、听骨链损伤。

\textbf{噪声对听觉外系统的影响：}
\begin{itemize}[leftmargin=*]
    \item \imp{神经系统}：头痛、头晕、失眠、多梦、记忆力减退
    \item \imp{心血管系统}：心率增快或减慢、血压升高（高血压患病率↑）
    \item \imp{消化系统}：胃肠功能紊乱、胃溃疡发病率↑
    \item \imp{内分泌系统}：ACTH、TSH等分泌异常
    \item \imp{生殖系统}：女性月经异常、妊娠不良结局↑
    \item \imp{心理影响}：烦躁、焦虑、注意力不集中、工作效率↓
\end{itemize}
\end{keybox}

\subsection{影响噪声危害的因素}

\begin{enumerate}[leftmargin=*]
    \item 噪声的强度和频谱特性（强度越大、频率越高，危害越大）
    \item 接触时间和方式（连续接触>间断接触）
    \item 噪声的类型（脉冲噪声危害>稳态噪声）
    \item 个体敏感性与个体差异
    \item 振动、高温、毒物等的联合作用
\end{enumerate}

\subsection{防止噪声危害的措施}

\begin{enumerate}[leftmargin=*]
    \item 控制和消除噪声源（技术革新、低噪设备）
    \item 控制噪声的传播和反射（隔声、吸声、消声、减振）
    \item 个体防护（耳塞、耳罩、防声帽）
    \item 健康监护（就业前体检、定期听力检查）
    \item 合理安排劳动和休息
    \item 制定和遵守噪声卫生标准（85dB(A)）
\end{enumerate}

% ----- 5.4 振动 -----
\section{生产性振动}

\subsection{基本概念}

\begin{defbox}
\textbf{振动（Vibration）}：质点或物体相对于平衡位置作的往复运动。

\textbf{振动参数*：}频率（Hz）、位移（mm）、速度（m/s）、加速度（m/s$^2$）。其中\keyword{加速度}与人体健康效应关系最密切。振动的频率决定了其对机体作用的部位——低频振动主要引起全身振动效应，高频振动主要引起局部振动效应。

\textbf{振动的分类：}
\begin{itemize}[leftmargin=*]
    \item \imp{全身振动}：通过身体支持面（站立的足、坐着的臀部）传导至全身
    \item \imp{局部振动（手传振动）}：通过手或手指作用于人体的振动
\end{itemize}
\end{defbox}

\subsection{局部振动病（手臂振动病）}

\begin{keybox}
\textbf{概念：}长期使用振动工具引起以\keyword{手部末梢循环障碍、末梢神经功能障碍和骨关节损害}为主要表现的疾病。

\textbf{主要接触机会：}凿岩工、油锯工、链锯工、砂轮磨光工等。

\textbf{发病机制*：}
\begin{enumerate}[leftmargin=*]
    \item \imp{末梢血管损伤}：振动刺激血管平滑肌→血管痉挛→血管内皮细胞损伤→血管壁增厚、管腔狭窄→末梢循环障碍
    \item \imp{末梢神经损伤}：振动直接损伤神经纤维和神经末梢→脱髓鞘改变→感觉传导速度减慢
    \item \imp{骨关节损伤}：振动使关节软骨和骨组织反复受力→微小损伤累积→关节退行性改变
    \item 寒冷是振动危害的重要\imp{诱发和加重因素}——寒冷使末梢血管收缩，加重缺血
\end{enumerate}

\textbf{手传振动的测量与评价（GBZ/T 189.9）：}
\begin{itemize}[leftmargin=*]
    \item 对多轴向振动，需分别测量X、Y、Z三个轴向的振动加速度
    \item \keyword{以三个轴向中最大振动加速度（最大值）作为评价依据}
    \item 评价指标：频率计权振动加速度有效值（a$_{hw}$）
    \item 日暴露时间、累积暴露时间直接影响危害程度
\end{itemize}

\textbf{影响振动危害的因素：}
\begin{itemize}[leftmargin=*]
    \item 振动频率和强度（高频振动危害更大）
    \item 接触振动的时间和方式
    \item 环境温度（寒冷加重振动危害——\keyword{振动性白指}在寒冷环境下易诱发）
    \item 机体状况
\end{itemize}

\textbf{临床表现——三大特征：}
\begin{enumerate}[leftmargin=*]
    \item \imp{末梢循环障碍}：\keyword{振动性白指（VWF）}——手指在寒冷刺激下出现发白（雷诺现象），是局部振动病最特征性改变，以食指、中指多见。发作顺序：苍白→青紫→潮红（反应性充血）
    \item \imp{末梢神经功能障碍}：手麻、触觉和痛觉减退、振动觉阈值升高
    \item \imp{骨关节和肌肉损害}：指关节变形、骨质疏松、骨质增生
\end{enumerate}

\textbf{诊断（GBZ 7-2014）：}明确的职业接触史+振动性白指+末梢神经功能障碍和/或骨关节损害。

\textbf{预防：}
\begin{enumerate}[leftmargin=*]
    \item 改革工艺和设备（减少手持振动工具的使用）
    \item 缩短接触时间（轮换作业）
    \item 改善环境温度（保暖、防寒）
    \item 个人防护（防振手套）
    \item 健康监护和就业禁忌证
\end{enumerate}
\end{keybox}

% ----- 5.5 辐射 -----
\section{非电离辐射与电离辐射}

\subsection{非电离辐射}

\begin{table}[H]
\centering
\caption{非电离辐射类型及健康效应}
\begin{tabularx}{\textwidth}{lX}
\hline
\textbf{类型} & \textbf{主要健康效应} \\
\hline
\imp{紫外线（UV）} & \textbf{电光性眼炎（急性角膜结膜炎）}——潜伏期4--12h，异物感、畏光、流泪；皮肤红斑和光感性皮炎；皮肤癌 \\
\imp{红外线} & 红外线白内障（玻璃工人白内障）；皮肤热损伤 \\
\imp{可见光} & 强光对视网膜的损伤 \\
\imp{激光} & 视网膜灼伤；角膜损伤 \\
\imp{射频辐射（微波）} & 热效应（组织加热）；非热效应（神经系统功能改变——神经衰弱综合征）\\
\imp{极低频电磁场} & 神经系统功能改变，潜在致癌风险（儿童白血病）\\
\hline
\end{tabularx}
\end{table}

\begin{defbox}
\textbf{射频电磁场的分类（按频率和波长）：}
\begin{table}[H]
\centering
\caption{射频电磁场分类}
\begin{tabularx}{\textwidth}{llll}
\hline
\textbf{频段名称} & \textbf{频率范围} & \textbf{波长范围} & \textbf{主要应用} \\
\hline
高频（HF） & 0.1--30 MHz & 3km--10m & 感应加热、短波通讯、广播 \\
超高频（VHF/UHF）& 30--300 MHz & 10m--1m & 电视广播、移动通讯 \\
微波（MW） & 300MHz--300GHz & 1m--1mm & 雷达、微波炉、卫星通讯 \\
\hline
\end{tabularx}
\end{table}

\textbf{射频辐射对机体的主要影响：}
\begin{itemize}[leftmargin=*]
    \item \imp{热效应}：射频能量被组织吸收转化为热能→局部或全身温度升高
    \item \imp{非热效应}：在不足以引起明显体温升高的低强度暴露下，仍可出现神经衰弱综合征（头痛、头晕、失眠、记忆力减退）、自主神经功能紊乱
    \item 敏感器官：眼（晶状体——微波白内障）、睾丸（精子生成障碍）
\end{itemize}
\end{defbox}

\subsection{电离辐射}

\begin{enumerate}[leftmargin=*]
    \item 外照射急性放射病（骨髓型、胃肠型、脑型）
    \item 外照射慢性放射病（以造血系统损害为主）
    \item 放射性皮肤损伤
    \item 远后效应：致癌（白血病、甲状腺癌、肺癌等）、遗传效应
\end{enumerate}

\textbf{防护三原则：}时间防护（缩短接触时间）、距离防护（增大与源距离）、屏蔽防护（铅板、混凝土等屏蔽材料）。

% ----- 5.6 复习题 -----
\section{复习题}

\begin{enumerate}[leftmargin=*, itemsep=8pt]
    \item \textbf{【名词解释】}高温作业；中暑；热射病；减压病；热适应；热习服
    \item \textbf{【名词解释】}噪声聋；振动性白指（VWF）；手臂振动病
    \item \textbf{【名词解释】}A计权声级（A声级）；等效连续A声级；听觉疲劳
    \item \textbf{【简答题】}简述高温作业的类型及高温对机体的主要生理影响。
    \item \textbf{【简答题】}简述中暑的四种临床分型及其特征（热射病、日射病、热痉挛、热衰竭）。
    \item \textbf{【简答题】}简述减压病的发病机制、临床表现和治疗原则。
    \item \textbf{【简答题】}简述噪声对听觉系统影响的进程（从听觉适应到噪声聋）。
    \item \textbf{【简答题】}简述手臂振动病的三大临床特征及发病机制。
    \item \textbf{【简答题】}简述防止噪声危害的主要措施。
    \item \textbf{【简答题】}简述电离辐射防护的三原则。
    \item \textbf{【简答题】}简述射频电磁场的分类及其健康效应。
\end{enumerate}

\subsection*{参考答案要点}

\begin{enumerate}[leftmargin=*, itemsep=5pt]
    \item \textbf{高温作业}：WBGT指数超过规定限值的作业。\textbf{中暑}：高温环境下热平衡和/或水盐代谢紊乱引起的急性疾病。\textbf{热射病}：最严重类型，高热>40$^\circ$C、无汗、意识障碍。\textbf{减压病}：高气压环境减压过快使溶解的氮气释放形成气泡栓塞引起的疾病。\textbf{热适应}：机体在长期高温环境下产生的形态、结构和生理功能的适应性变化。\textbf{热习服}：短期反复接触高温后生理功能补偿性调整（7--14天），脱离高温后消退。
    \item \textbf{噪声聋}：长期接触生产性噪声引起的以高频听力下降为主的感音性耳聋。\textbf{VWF}：手指在寒冷刺激下出现的苍白现象（雷诺现象），局部振动病最特征性改变。\textbf{手臂振动病}：长期使用振动工具引起的手部末梢循环障碍、末梢神经功能障碍和骨关节损害。
    \item \textbf{A声级}：用A计权网络测量的声压级，模拟人耳对40方纯音的响应，dB(A)。\textbf{等效连续A声级}：能量等效于实际噪声的连续稳态A声级。\textbf{听觉疲劳}：长时间暴露于强噪声后听力明显下降，脱离后数小时至十几小时恢复。
    \item 三类：干热型（高温、强辐射、低湿）、湿热型（高温、高湿、辐射不强）、夏季露天作业。影响：体温调节（出汗蒸发散热为主）、循环（心率↑、血液重新分配）、消化（分泌↓）、泌尿（浓缩）、神经（抑制→作业能力↓）。
    \item 热射病：体温>40$^\circ$C、汗闭、昏迷、CNS症状（最严重，死亡率高）。日射病：夏季露天作业、头（颅）温升高、有汗、恶心呕吐、昏迷。热痉挛：大量出汗补水未补盐→低钠→肌肉痉挛（腓肠肌多见）、体温不高、无昏迷。热衰竭：心血管功能失代偿→冷汗、体温↑、昏厥。
    \item 机制：高压下氮气大量溶解→减压过快→氮气形成气泡→栓塞血管和压迫组织→缺血缺氧。表现：皮肤瘙痒、屈肢症（关节肌肉剧痛）、脊髓损害。治疗：加压舱加压治疗（唯一有效）。
    \item 进程：听觉适应（生理性，数分钟恢复）→听觉疲劳（听阈提高15--30dB，数小时恢复）→永久性听阈位移（NIPTS，不可逆）→噪声聋（先高频↓，4kHz听谷为早期特征，后语言频率↓）。
    \item 三大特征：末梢循环障碍（振动性白指VWF——寒冷诱发手指苍白→青紫→潮红）；末梢神经功能障碍（手麻、触觉振动觉减退）；骨关节肌肉损害（指关节变形、骨质疏松）。
    \item 控制和消除噪声源；控制噪声传播（隔声、吸声、消声、减振）；个体防护（耳塞、耳罩）；健康监护（定期听力检查）；合理安排劳动休息；执行卫生标准（85dB(A)）。
    \item 时间防护（缩短接触时间）、距离防护（增大与源距离）、屏蔽防护（铅板、混凝土等屏蔽材料）。
\end{enumerate}

\newpage
% =====================================================================
%  CHAPTER 6: 职业性致癌因素与职业性肿瘤
% =====================================================================
\newpage